\documentclass[10pt,a4paper]{amsart}
\usepackage[margin=19mm]{geometry}
\usepackage{amsmath,amssymb,mathtools}
\usepackage[scr=rsfs,cal=euler]{mathalfa}
\usepackage{xfrac}
\usepackage[unicode,hidelinks,bookmarks=false]{hyperref}
\newcommand{\ie}{i.e.\ }
\newcommand{\cf}{cf.\ }
\newcommand{\resp}{respectively\ }
\newcommand{\Subsection}{\subsection}
\providecommand{\nobreakdash}{\nobreak\hbox{-}}
\setlength{\emergencystretch}{2em}
\multlinegap0pt
\usepackage{bm}
\usepackage{bbm}
\usepackage{dsfont}
\usepackage{xspace}
\usepackage{cancel}
\usepackage{mathtools}
\usepackage{amsthm}
\makeatletter
\@ifundefined{theorem}{\newtheorem{theorem}{Theorem}}{}
\@ifundefined{lemma}{\newtheorem{lemma}[theorem]{Lemma}}{}
\@ifundefined{prop}{\newtheorem{prop}[theorem]{Proposition}}{}
\makeatother
\newcommand{\E}{\mathbb{E}}
\newcommand{\Tf}{T}
\renewcommand{\P}{\mathbb{P}}
\title[Random sequential nearest-neighbor coloring on trees]{Random sequential nearest-neighbor coloring on trees}
\author{Anne-Laure Basdevant, Arvind Singh}
\date{Source: arXiv:2606.24793v1 (CC BY 4.0)}
\begin{document}
\maketitle

\noindent\emph{Source and licence.} Random sequential nearest-neighbor coloring on trees. Version arXiv:2606.24793v1 (CC BY 4.0), distributed under the Creative Commons Attribution 4.0 International licence, as recorded in the arXiv licence metadata for this exact version and shown on the abstract page https://arxiv.org/abs/2606.24793v1. Full rights notice: licenses/arxiv-2606.24793-CC-BY-4.0.txt.\par\smallskip

\noindent\textit{Editorial scope.} Counted unit: the Prop whose source label is \texttt{prop:unsommetvide} in the article (source0.tex lines 1657--1661, source label \texttt{prop:unsommetvide}), delivered together with the article's own complete original proof of it (source0.tex lines 1665--1905, one \texttt{proof} environment of 241 lines). No mathematics is written, repaired or re-derived: statement and proof are the article's own text line for line. The proof is detached from the statement in the source: the article states the result at lines 1657--1661 and prints its proof at lines 1665--1905, and the proof environment's own header names the statement, so the association is the article's own. Required context retained uncounted: eq:defai (lines 1528--1538). Every definition, display and label of those lines is retained unchanged. Omitted from this package and not redistributed: 1--1527, 1539--1656, 1662--1664, 1906--2041. Nothing inside a retained excerpt is omitted or altered: every retained body is delivered exactly as the article's lines are written, so no omission receipt of any kind accompanies this package. Citations: the retained excerpts carry no citation command at all, so no bibliography entry of the article is transcribed and no reference is redistributed. Reference adaptation: none was needed and none was made; every reference inside the delivered unit resolves inside it, so no unresolved reference or citation is produced. Printed slips kept literal: no printed word, symbol, hypothesis or number of the counted statement or of its proof was altered. Layout and class: the article's own class is \texttt{article} with packages including inputenc, fontenc, lmodern, geometry, babel, fullpage, amssymb, amsmath, mathtools, amsthm, dsfont, stmaryrd, bm, graphicx; the wrapper is the lane's pinned \texttt{amsart} editorial wrapper at 10pt on A4, which restates the theorem-like environments the retained text needs and adds the article's own macro definitions that the retained text needs (\texttt{\textbackslash E}, \texttt{\textbackslash P}, \texttt{\textbackslash Tf}), with extra wrapper packages (bm, bbm, dsfont, xspace, cancel, mathtools). The delivered numbering is the unit's own. Assets: no retained span contains a figure environment, a graphics-inclusion command, a PDF-page inclusion, a table or a TikZ picture; no image file is copied into this package, the package declares no asset dependency and no image file is redistributed with it. Licence: version arXiv:2606.24793v1 of this article is distributed under the Creative Commons Attribution 4.0 International licence, as recorded in the arXiv licence metadata for this exact version and shown on the abstract page https://arxiv.org/abs/2606.24793v1. A full rights notice is delivered at licenses/arxiv-2606.24793-CC-BY-4.0.txt. Characters: every retained excerpt is pure ASCII, so the delivered engine, XeLaTeX with its default Latin Modern text font, reports no missing character.
\begin{equation}\label{eq:defai}
\begin{aligned}
\mathcal{A}_i := \Bigl\{&
X^{\Tf}_i \text{ is colored blue},\\
& |\Tf|/2 - H < |X^{\Tf}_i| < |\Tf|/2,\\
& \text{at time }(\tau_{X^{\Tf}_i})^{-},\ 
\Tf_{X^{\Tf}_i} \text{ contains at most } C \text{ colored non-leaf vertices,}\\
& \text{and all these vertices are at distance at most } 2H \text{ from the leaves}
\Bigr\}.
\end{aligned}
\end{equation}

\begin{prop}\label{prop:unsommetvide} Let $\eta>0$. Then, there exist $H,C,N,n_0\ge 0$ such that, if $|\Tf|\ge n_0$, then for any choice of vertices $w_1,\ldots,w_C$ of height greater than or equal to $|\Tf|-2H$, starting from $B_0=\{o\}$ and $R_0=\partial \Tf \cup\{w_1,\ldots,w_C\}$, we have
\begin{equation*}
 \P(\forall 1\le i\le N,\; \mathcal{A}_i^c)\le \eta.
\end{equation*}
\end{prop}

\begin{proof}[Proof of Proposition \ref{prop:unsommetvide}]

For $v\in \Tf$ and $r,h\ge 0$, we denote
$$
B(v,r):=\{w\in \Tf:\ d(v,w)\le r\}
\qquad \text{and} \qquad
B_{\ge h}(v,r):=\{w\in \Tf:\ d(v,w)\le r \ \text{and}\ |w|\ge h\},
$$
the ball centered at $v$ of radius $r$ and the set of vertices in this ball whose height is at least $h$, respectively.
For $v_1,v_2\in \Tf$, we denote by $v_1\wedge v_2$ their most recent common ancestor.
Recall that $R_0=\partial \Tf \cup\{w_1,\ldots,w_C\}$ denotes the set of initial red vertices, and that $X^{\Tf}_1,X^{\Tf}_2,\ldots$ enumerates the vertices of $\Tf\setminus\{o\}$ whose heights are strictly smaller than $|\Tf|/2$, ordered according to their coloring times. Since $\Tf$ is fixed throughout this section, we will simplify the notation and write $X_1,X_2,\ldots$ instead of $X^{\Tf}_1,X^{\Tf}_2,\ldots$.
We now introduce the following events.



\begin{eqnarray*}
\mathcal{E}_1=\mathcal{E}_1(\Tf,N,H,C,R_0) &:=& \{\forall i\neq j \le N, |X_i\wedge X_j| < |\Tf|/8\},\\
\mathcal{E}_2=\mathcal{E}_2(\Tf,N,H,C,R_0) &:=& \{\forall i \le N, |X_i|> |\Tf|/2-H\},\\
\mathcal{E}_3=\mathcal{E}_3(\Tf,N,H,C,R_0) &:=& \{\forall i \le N, \forall k \le C, X_i\notin B(w_k,|\Tf|/2+H)\},\\
\mathcal{E}_4=\mathcal{E}_4(\Tf,N,H,C,R_0) &:=&\{\forall i \le N, \mathcal{A}_i^c\}\cap  \mathcal{E}_1 \cap  \mathcal{E}_2 \cap  \mathcal{E}_3.
\end{eqnarray*}
Thus, we have
$$
\P\bigl(\forall i \le N,\; \mathcal{A}_i^c\bigr)
\le \P(\mathcal{E}_1^c)+\P(\mathcal{E}_2^c)+\P(\mathcal{E}_3^c)+\P(\mathcal{E}_4).
$$
The first three terms on the right-hand side are easy to bound.

Indeed, to bound $\P(\mathcal{E}_1^c)$, observe that $X_i$ and $X_j$ are distributed as two distinct vertices chosen uniformly at random from the complete $K$-ary tree of height $\lceil |\Tf|/2-1 \rceil$, excluding the root. It is easy to see that the height of the most recent common ancestor of such vertices is stochastically dominated by a geometric random variable with parameter $(K-1)/K$. Hence, we obtain
\begin{equation}\label{eq:E1}
\P(\mathcal{E}_1^c)
\le N^2 \P\bigl( |X_1 \wedge X_2| \ge |\Tf|/8 \bigr)
\le N^2 \frac{1}{K^{|\Tf|/8}}.
\end{equation}

Concerning $\P(\mathcal{E}_2^c)$, and still using that $X_i$ is uniformly distributed on the complete $K$-ary tree of height $\lceil |\Tf|/2-1 \rceil$, excluding the root, together with a simple counting argument on the number of vertices of height smaller than $|\Tf|/2-H$, we obtain
\begin{equation}\label{eq:E2}
\P(\mathcal{E}_2^c)\le N K^{-H}.
\end{equation}

For $\P(\mathcal{E}_3^c)$, recall that by assumption, for any $k\le C$, $|w_k|\ge |\Tf|-2H$, whereas for any $i\ge 1$ we have $|X_i|\le \lceil |\Tf|/2-1 \rceil$. Let $z_k$ denote the ancestor of $w_k$ at height $\lceil |\Tf|/2-1 \rceil$. Then
$$
X_i\in B(w_k,|\Tf|/2+H)\ \Longrightarrow\ X_i\in B(z_k,3H).
$$
Hence,
\begin{equation}\label{eq:E3}
\P(\mathcal{E}_3^c)
\le CN\,\P\bigl(X_1\in B(z_1,3H)\bigr)
\le CN \frac{K^{3H}}{K^{|\Tf|/2-2}},
\end{equation}
again using that $X_1$ is uniformly distributed on the complete $K$-ary tree of height $\lceil |\Tf|/2-1 \rceil$ (excluding the root).
\medskip


It remains to upper bound $\P(\mathcal{E}_4)$, which requires more work. In the following, for a vertex $v$ and a time $t>0$, we say that the clock of $v$ has rung before time $t$ if $0<\tau_v<t$. We introduce the following two events:
\begin{itemize}
\item[(i)] $\mathcal{F}_i := \{\text{At time } (\tau_{X_i})^{-}, \text{ strictly more than } C \text{ clocks have rung in } \Tf_{X_i}\}$.
\item[(ii)] $\mathcal{O}_i := \{\text{A clock has rung in } B_{\ge |\Tf|/2}(X_i, |X_i|) \text{ before time } (\tau_{X_i})^{-}\}$.
\end{itemize}



\begin{lemma} Assume that $H< |\Tf|/4$. Then
   we have $$\mathcal{E}_4\subset \{\forall i\le N, \mathcal{O}_i \cup \mathcal{F}_i\}\cap \mathcal{E}_1\cap \mathcal{E}_2.$$
\end{lemma}


\begin{proof}

Recalling the definition of $\mathcal{A}_i$ given in \eqref{eq:defai}, we have
$
\mathcal{A}_i := \mathcal{X}_i \cap \mathcal{Y}_i \cap \mathcal{Z}_i,
$
where
\begin{eqnarray*}
\mathcal{X}_i &:=& \{|X_i|>|\Tf|/2-H\},\\
\mathcal{Y}_i &:=& \{\text{$X_i$ is colored blue}\},\\
\mathcal{Z}_i &:=& \Bigl\{\text{At time } (\tau_{X_i})^{-},\; \Tf_{X_i} \text{ contains at most } C \text{ colored vertices that are not leaves,}\\
&& \qquad\quad \text{and each of them is at distance at most } 2H \text{ from the leaves}\Bigr\}.
\end{eqnarray*}
We prove that
$$
\mathcal{F}_i^c \cap \mathcal{O}_i^c \cap \mathcal{E}_1 \cap \mathcal{E}_2 \cap \mathcal{E}_3 \subset \mathcal{A}_i,
$$
by showing that the event on the left-hand side implies $\mathcal{X}_i$, $\mathcal{Y}_i$, and $\mathcal{Z}_i$.

\medskip

\noindent We have $\mathcal{E}_2 \subset \mathcal{X}_i$. Assume now that $\mathcal{O}_i^c \cap \mathcal{E}_1 \cap \mathcal{E}_2 \cap \mathcal{E}_3$ occurs. Observe that for $i \neq j \le N$,
$$
d(X_i,X_j)
= |X_i| + |X_j| - 2|X_i \wedge X_j|
> |X_i| + \frac{|\Tf|}{2} - H - \frac{|\Tf|}{4}
> |X_i|.
$$
Hence, we have $X_j \notin B(X_i,|X_i|)$. By definition of the sequence $(X_j)_{j\le N}$, it follows that at time $(\tau_{X_i})^{-}$ no clock in $B(X_i,|X_i|)$ at height smaller than $|\Tf|/2$ has rung. Moreover, since $\mathcal{O}_i^c$ holds, no clock in $B(X_i,|X_i|)$ at height at least $|\Tf|/2$ has rung either. Therefore, no clock in $B(X_i,|X_i|)$ has rung before time $(\tau_{X_i})^{-}$. Using now $\mathcal{E}_3$ and the fact that $|X_i|<|\Tf|/2$, we deduce that $B(X_i,|X_i|)$ contains no vertex of $R_0$. Hence, at time $(\tau_{X_i})^{-}$, the ball $B(X_i,|X_i|)$ contains no red vertex but does contain the root, which is blue. Consequently, $X_i$ becomes blue, and we obtain
$$
\mathcal{O}_i^c \cap \mathcal{E}_1 \cap \mathcal{E}_2 \cap \mathcal{E}_3 \subset \mathcal{Y}_i.
$$

Moreover, on the event $\mathcal{O}_i^c \cap \mathcal{E}_1 \cap \mathcal{E}_2$, we have $|X_i|>|\Tf|/2-H$ and no clock has rung in $B(X_i,|X_i|)$ before time $\tau_{X_i}$. In particular, any clock that has rung in $\Tf_{X_i}$ before time $\tau_{X_i}$ must correspond to a vertex at height at least $2|X_i|\ge |\Tf|-2H$.
Hence, by definition of $\mathcal{F}_i$, on the event
$\mathcal{F}_i^c \cap \mathcal{O}_i^c \cap \mathcal{E}_1 \cap \mathcal{E}_2$,
at most $C$ vertices have rung in $\Tf_{X_i}$ before time $\tau_{X_i}$, and all of them lie at distance at most $2H$ from the leaves. Using additionally $\mathcal{E}_3$, which ensures that the initial red vertices in $\Tf_{X_i}$ are only leaves, we obtain
$$
\mathcal{F}_i^c \cap \mathcal{O}_i^c \cap \mathcal{E}_1 \cap \mathcal{E}_2 \cap \mathcal{E}_3
\subset \mathcal{Z}_i.
$$
Combining the previous observations, we deduce that
$$
\mathcal{F}_i^c \cap \mathcal{O}_i^c \cap \mathcal{E}_1 \cap \mathcal{E}_2 \cap \mathcal{E}_3
\subset \mathcal{A}_i,
$$
which implies
$$
\mathcal{A}_i^c \cap \mathcal{E}_1 \cap \mathcal{E}_2 \cap \mathcal{E}_3
\subset \mathcal{F}_i \cup \mathcal{O}_i.
$$
Hence,
$$
\mathcal{E}_4 \subset \Bigl\{\forall i \le N,\; \mathcal{F}_i \cup \mathcal{O}_i\Bigr\} \cap \mathcal{E}_1 \cap \mathcal{E}_2.
$$
This concludes the proof of the lemma.
\end{proof}

Let us now derive an upper bound on the probability of
$\{\forall i \le N,\; \mathcal{F}_i \cup \mathcal{O}_i\}\cap \mathcal{E}_1\cap \mathcal{E}_2$.
Recall that $\tau_v$ denotes the coloring time of a vertex $v$. We define the $\sigma$-algebra
$$
\mathcal{H}
:= \sigma\bigl(\{\tau_v:\ |v|<|\Tf|/2\}\bigr),
$$
so that $\mathcal{E}_1,\mathcal{E}_2,\mathcal{E}_3 \in \mathcal{H}$. For $i\le N$, we denote by
$\tau_i$ the coloring time of $X_i$, i.e. $\tau_i=\tau_{X_i}$. Note that $\tau_i$ is $\mathcal{H}$-measurable.
We observe that for any $i\le N$ and any $u\in B_{\ge |\Tf|/2}(X_i,|X_i|)$, we have
$$
2|u\wedge X_i|
= |u| + |X_i| - d(u,X_i)
\ge |u|
\ge \frac{|\Tf|}{2}.
$$
Moreover, if $u\in \Tf_{X_i}$, then $|u\wedge X_i|=|X_i|$. Hence, assuming $H<|\Tf|/4$, on the event $\mathcal{E}_2$ we obtain that for all
$u\in B_{\ge |\Tf|/2}(X_i,|X_i|)\cup \Tf_{X_i}$,
$$
|u\wedge X_i|>\frac{|\Tf|}{8}.
$$
In particular, on $\mathcal{E}_1 \cap \mathcal{E}_2$, for any $i\neq j \le N$, the sets
$B_{\ge |\Tf|/2}(X_i,|X_i|)\cup \Tf_{X_i}$ and
$B_{\ge |\Tf|/2}(X_j,|X_j|)\cup \Tf_{X_j}$
are disjoint. Therefore, conditionally on $\mathcal{H}$ and on $\mathcal{E}_1 \cap \mathcal{E}_2$, the events $\mathcal{F}_i \cup \mathcal{O}_i$, $i\le N$, are independent. Besides, we have
$$
\sharp B_{\ge |\Tf|/2}(X_i, |X_i|) \le (K+1)K^{|X_i|}
\qquad \text{and} \qquad
\sharp\Tf_{X_i} \le K^{|\Tf|/2+H+1}.
$$
Hence, conditionally on $\mathcal{H}$ and on the event $\mathcal{E}_1 \cap \mathcal{E}_2$, the time until the first clock rings in $B_{\ge |\Tf|/2}(X_i, |X_i|)$,
namely
\[
\min\{\tau_v:\ v\in B_{\ge |\Tf|/2}(X_i, |X_i|)\},
\]
is stochastically lower bounded by an exponential random variable with parameter $(K+1)K^{|X_i|}$. Moreover, the number of clocks that have rung before time $\tau_i$ in $\Tf_{X_i}$ is stochastically upper bounded by a binomial random variable with parameters $\bigl(K^{|\Tf|/2+H+1},\, 1-e^{-\tau_i}\bigr)$. Let $e_h$ denote an exponential random variable with parameter $(K+1)K^{h}$ and let $Y_t$ be a binomial random variable with parameters $\bigl(K^{|\Tf|/2+H+1},\, 1-e^{-t}\bigr)$. We obtain
\begin{eqnarray*}
\P(\mathcal{E}_4\mid \mathcal{H})&\le& 1_{\{\mathcal{E}_1\cap \mathcal{E}_2\}} \prod_{i=1}^N\P(e_{|X_i|}\le \tau_i \mbox{ or } Y_{\tau_i}\ge C \mid \mathcal{H})\\
&=& 1_{\{\mathcal{E}_1\cap \mathcal{E}_2\}} \prod_{i=1}^N\Big(1-\P(e_{|X_i|}> \tau_i \mid \mathcal{H})\P( Y_{\tau_i}<C \mid \tau_i)\Big).\end{eqnarray*}
We write, for a fixed $A>0$,
$$\P( Y_{\tau_i}<C \mid \tau_i)\ge 1_{\{\tau_i\le A K^{-|\Tf|/2}\}}\P(Y_{A K^{-|\Tf|/2}}<C).$$
We set $q_{A,C,|\Tf|}:=\P(Y_{A K^{-|\Tf|/2}}<C)$ so that 
\begin{eqnarray*}
\P(\mathcal{E}_4\mid \mathcal{H})&\le&  1_{\{\mathcal{E}_1\cap \mathcal{E}_2\}} \prod_{i=1}^N\left(1-\exp(-(K+1)K^{|X_i|}\tau_i)q_{A,C,|\Tf|}1_{\{\tau_i\le A K^{-|\Tf|/2}\}}\right).\end{eqnarray*}
Let $m$ denote the largest integer strictly smaller than $|\Tf|/2$, and let $\Tf_m^*$ be the set of vertices of $\Tf\setminus\{o\}$ whose height is at most $m$. We have
$$
\sharp\Tf_m^*=\frac{K(K^{m}-1)}{K-1}.
$$
Note that $X_1,\ldots,X_N$ are $N$ distinct vertices of $\Tf_m^*$ chosen uniformly at random. Moreover, the random variables $(\tau_i)_{i\le N}$ and $(X_i)_{i\le N}$ are independent, and $(\tau_i-\tau_{i-1})_{i\le N}$ are independent exponential random variables with parameters $\sharp\Tf_m^* - i$. We obtain, for any $i\ge 1$ and $h<m$,
$$
\P(|X_i|=m-h)
=
\frac{K^{m-h}}{\sharp\Tf_m^*}
=
\frac{(K-1)K^{m-h}}{K(K^m-1)}
\ge \frac{1}{2K^h}.
$$
Moreover, for any $h_i\ge 1$, we have
\begin{eqnarray*}
\P((|X_1|,\ldots,|X_N|)=(h_1,\ldots,h_N))&\le& \left(\frac{\sharp\Tf_m^*}{\sharp\Tf_m^* - N}\right)^N\prod_{i=1}^N\P(|X_i|=h_i)\\
&\le & \left(1-\frac{N}{K^{m}}\right)^{-N}\prod_{i=1}^N\P(|X_i|=h_i)\\
&\le & \left(1-\frac{N}{K^{|\Tf|/2-1}}\right)^{-N}\prod_{i=1}^N\P(|X_i|=h_i).
\end{eqnarray*}
Thus, we get 
\begin{eqnarray*}
\P(\mathcal{E}_4\mid (\tau_i)_{i\le N})&\le& \sum_{h_1,\ldots,h_N\ge 1} \P((|X_1|,\ldots,|X_N|)=(h_1,\ldots,h_N))\prod_{i=1}^N\left (1-e^{-\tau_i(K+1)K^{h_i}}q_{A,C,|\Tf|}1_{\{\tau_i\le A K^{-|\Tf|/2}\}}\right)\\
&\le &\left(1-\frac{N}{K^{|\Tf|/2-1}}\right)^{-N}\sum_{h_1,\ldots,h_N\ge 0} \prod_{i=1}^N\left (\P(|X_i|=h_i)(1-e^{-\tau_i(K+1)K^{h_i}}q_{A,C,|\Tf|}1_{\{\tau_i\le A K^{-|\Tf|/2}\}})\right)\\
&\le& \left(1-\frac{N}{K^{|\Tf|/2-1}}\right)^{-N}\prod_{i=1}^N \E\left(1-e^{-\tau_i(K+1)K^{|X_i|}}q_{A,C,|\Tf|}1_{\{\tau_i\le A K^{-|\Tf|/2}\}} \mid (\tau_i)_{i\le N}
\right)\\
&=&  \left(1-\frac{N}{K^{|\Tf|/2-1}}\right)^{-N}\prod_{i=1}^N \left(1-\E\left(e^{-\tau_i(K+1)K^{|X_i|}} \mid (\tau_i)_{i\le N}\right)q_{A,C,|\Tf|}1_{\{\tau_i\le A K^{-|\Tf|/2}\}}
\right)\\
&\le&  \left(1-\frac{N}{K^{|\Tf|/2-1}}\right)^{-N}\prod_{i=1}^N \left(1-\left(\sum_{h=0}^{L}\frac{1}{2K^h}e^{-\tau_i(K+1)K^{m-h}}\right) q_{A,C,|\Tf|}1_{\{\tau_i\le A K^{-|\Tf|/2}\}}
\right)
\end{eqnarray*}
valid for any $L<m$ (recall that $m$ is the largest integer strictly smaller than $|\Tf|/2$).
\medskip

We now use that the random variables $(\tau_i-\tau_{i-1})_{i\le N}$ are independent exponential random variables with parameters $\sharp\Tf_m^* - i$. If $i\le \sharp\Tf_m^*/2$, we have
$$\sharp\Tf_m^* - i\ge  \frac{\sharp\Tf_m^*}{2}\ge K^{m-1}\ge K^{|\Tf|/2-2}.$$
Thus, if $N\le K^m/2$, the sequence $(K^{|\Tf|/2-2}(\tau_i-\tau_{i-1}))_{i\le N}$ is stochastically smaller than a sequence $(\varepsilon_i)_{i\le N}$ of i.i.d.\ exponential random variables with parameter 1. Let $\sigma_i=\sum_{k=1}^i\varepsilon_k$.
We get 
\begin{eqnarray*}
\P(\mathcal{E}_4)
&\le&  \left(1-\frac{N}{K^{|\Tf|/2-1}}\right)^{-N}\E\left(\prod_{i=1}^N \left(1-\left(\sum_{h=0}^L\frac{1}{2K^h}e^{-\tau_i(K+1)K^{|\Tf|/2-h}}\right) q_{A,C,|\Tf|}1_{\{\tau_i\le A K^{-|\Tf|/2}\}}
\right)\right)\\
&\le& \left(1-\frac{N}{K^{|\Tf|/2-1}}\right)^{-N}\E\left(\prod_{i=1}^N \left(1-\left(\sum_{h=0}^L\frac{1}{2K^h}e^{-\sigma_i(K+1)K^{2-h}}\right) q_{A,C,|\Tf|}1_{\{\sigma_i\le A K^{-2}\}}
\right)\right).
\end{eqnarray*}
We now use the following lemma, whose proof is postponed until the end of this section.
\begin{lemma}\label{lem:technique}
    Let $\eta>0$. Then, there exist integers $N,L>0$ such that if $(\varepsilon_i)_{i\le N}$ are i.i.d.\ exponential random variables with parameter 1 and $\sigma_i=\sum_{k=1}^i\varepsilon_k$, we have 
    \begin{equation}\label{eq:technique}
        \E\left(\prod_{i=1}^N \left(1-\sum_{h=0}^L\frac{1}{2K^h}e^{-\sigma_i(K+1)K^{2-h}}\right) \right)\le \eta/5.
    \end{equation}    
\end{lemma}
We finish the proof of Proposition \ref{prop:unsommetvide} assuming that Lemma \ref{lem:technique} holds. Fix $\eta>0$ and choose $N,L$ such that \eqref{eq:technique} holds. Since for each $i\le N$, $1_{\{\sigma_i\le A K^{-2}\}}$ tends a.s. to 1 as $A$ tends to infinity, by the dominated convergence theorem, there exists $A$ large enough such that 
\begin{equation*}
        \E\left(\prod_{i=1}^N \left(1-1_{\{\sigma_i\le A K^{-2}\}}\sum_{h=0}^L\frac{1}{2K^h}e^{-\sigma_i(K+1)K^{2-h}}\right) \right)\le \eta/4.
    \end{equation*}
Using \eqref{eq:E2}, fix now $H$ such that 
\begin{equation*}
\P(\mathcal{E}_2^c)\le N\frac{1}{K^{H}}\le \eta/3.
\end{equation*}
Recall that 
$$q_{A,C,|\Tf|}:=\P(Y_{A K^{-|\Tf|/2}}<C) \mbox{ where } Y_{A K^{-|\Tf|/2}}\sim \mbox{Bin}(K^{|\Tf|/2+H+1}, 1-e^{-A K^{-|\Tf|/2}}).$$
Thus, we get that
$$1-q_{A,C,|\Tf|}=\P(Y_{A K^{-|\Tf|/2}}\ge C)\le \frac{\E(Y_{A K^{-|\Tf|/2}})}{C}\le \frac{A K^{H+1}}{C}.$$
This shows that, for fixed $A$ and $H$, we can find $C$ large enough such that $q_{A,C,|\Tf|}$ is close to 1. Using again the dominated convergence theorem, we deduce the existence of a $C$ such that 
$$\E\left(\prod_{i=1}^N \left(1-\left(\sum_{h=0}^L\frac{1}{2K^h}e^{-\sigma_i(K+1)K^{2-h}}\right) q_{A,C,|\Tf|}1_{\{\sigma_i\le A K^{-2}\}}
\right)\right)\le \eta/3,$$
and so, for this choice of $N,H,C$, we obtain $\P(\mathcal{E}_4)\le \eta/2$ if $|\Tf|$ is large enough. Combining this with \eqref{eq:E1} and \eqref{eq:E3}, we conclude that for this choice of $H,C,N$ and for $|\Tf|$ large enough,
\begin{equation*}
 \P(\forall 1\le i\le N,\; \mathcal{A}_i^c)\le \eta.
\end{equation*}
This is exactly the statement of Proposition \ref{prop:unsommetvide}.
\end{proof}

\end{document}
